% ================================================================
%  sesion04.tex  —  Sesión 4 de 20
%  Coordenadas y localización
%  Almería en Cifras y Formas · Matemáticas 1.º ESO
%  Compilar con:  pdflatex sesion04.tex  (dos pasadas)
% ================================================================
\documentclass[aspectratio=169, 10pt, spanish]{beamer}
\usepackage{almeria-beamer}

% ---------------------------------------------------------------
%  DATOS DE LA SESIÓN
% ---------------------------------------------------------------
\renewcommand{\SessionNumber}{04}
\renewcommand{\SessionTitle}{Coordenadas y localización}
\renewcommand{\SessionName}{Sesión 04}
\renewcommand{\SessionObjective}{%
  Leer, escribir y localizar lugares exactos en un mapa de Almería
  usando pares de coordenadas cartesianas $(x, y)$.}
\renewcommand{\BlockName}{Bloque 1: Mapa y coordenadas}

% ================================================================
\begin{document}

% ---------------------------------------------------------------
%  PORTADA
% ---------------------------------------------------------------
\begin{frame}[plain]
  \titlepage
\end{frame}

% ---------------------------------------------------------------
%  ÍNDICE DE LA SESIÓN
% ---------------------------------------------------------------
\begin{frame}{Índice de la sesión}
  \begin{columns}[t]
    \begin{column}{0.5\textwidth}
      \begin{enumerate}
        \item El sistema de coordenadas cartesianas
        \item Leer un par $(x, y)$
        \item Los cuatro cuadrantes y sus signos
        \item Representar puntos en el plano
      \end{enumerate}
    \end{column}
    \begin{column}{0.5\textwidth}
      \begin{enumerate}
        \setcounter{enumi}{4}
        \item Escala: de la cuadrícula al mapa real
        \item Ejercicios multinivel (Bloom)
        \item Evidencia del día
        \item Error frecuente
        \item Resumen
      \end{enumerate}
    \end{column}
  \end{columns}
\end{frame}

% ---------------------------------------------------------------
\seccionframe{1. El sistema de coordenadas cartesianas}
% ---------------------------------------------------------------

\begin{frame}{Ejes cartesianos: eje X y eje Y}
  \begin{columns}[c]
    \begin{column}{0.50\textwidth}
      \begin{definicionbox}{Sistema de coordenadas cartesianas}
        Dos rectas perpendiculares que se cortan en el
        \textbf{origen} $(0,0)$:
        \begin{itemize}
          \item \textbf{Eje X} (horizontal): valores positivos
                a la derecha, negativos a la izquierda.
          \item \textbf{Eje Y} (vertical): valores positivos
                hacia arriba, negativos hacia abajo.
        \end{itemize}
      \end{definicionbox}
      \vspace{6pt}
      \begin{teoriabox}
        \AlmFontSmall
        \textbf{Regla mnemotécnica:}\\
        «Primero camina por la \textbf{calle} (eje X),
        luego sube o baja en el \textbf{ascensor} (eje Y).»
      \end{teoriabox}
    \end{column}
    \begin{column}{0.46\textwidth}
      \begin{center}
        \begin{tikzpicture}[scale=0.72]
          % Ejes
          \draw[recto, ->] (-3.4,0) -- (3.6,0)
            node[right, font=\AlmFontSmall\bfseries, color=AlmAzulM] {$x$};
          \draw[recto, ->] (0,-3.4) -- (0,3.6)
            node[above, font=\AlmFontSmall\bfseries, color=AlmAzulM] {$y$};
          % Cuadrícula ligera
          \foreach \i in {-3,-2,-1,1,2,3}{
            \draw[AlmAzulC!25, line width=0.3pt]
              (\i,-3.3) -- (\i,3.3);
            \draw[AlmAzulC!25, line width=0.3pt]
              (-3.3,\i) -- (3.3,\i);
          }
          % Marcas en los ejes
          \foreach \i in {-3,-2,-1,1,2,3}{
            \draw[AlmAzulM, line width=0.6pt]
              (\i,-0.08) -- (\i,0.08);
            \draw[AlmAzulM, line width=0.6pt]
              (-0.08,\i) -- (0.08,\i);
            \node[font=\AlmFontTiny, color=AlmGris] at (\i,-0.28) {\i};
            \node[font=\AlmFontTiny, color=AlmGris] at (-0.28,\i) {\i};
          }
          % Origen
          \node[font=\AlmFontTiny, color=AlmGris] at (-0.28,-0.28) {0};
          % Indicadores de dirección
          \node[font=\AlmFontTiny\bfseries, color=AlmTerra] at (2.2,-0.45)
            {$+$};
          \node[font=\AlmFontTiny\bfseries, color=AlmTerra] at (-2.2,-0.45)
            {$-$};
          \node[font=\AlmFontTiny\bfseries, color=AlmTerra] at (0.35,2.3)
            {$+$};
          \node[font=\AlmFontTiny\bfseries, color=AlmTerra] at (0.35,-2.3)
            {$-$};
        \end{tikzpicture}
      \end{center}
    \end{column}
  \end{columns}
\end{frame}

% ---------------------------------------------------------------
\seccionframe{2. Leer un par $(x, y)$}
% ---------------------------------------------------------------

\begin{frame}{¿Qué significa el par $(x, y)$?}
  \begin{columns}[c]
    \begin{column}{0.50\textwidth}
      \begin{definicionbox}{Par de coordenadas $(x,\,y)$}
        \begin{itemize}
          \item \textbf{Primera coordenada} $x$:\\
                desplazamiento \emph{horizontal}
                desde el origen.\\
                $(+)$ = derecha, $(-)$ = izquierda.
          \item \textbf{Segunda coordenada} $y$:\\
                desplazamiento \emph{vertical}
                desde el origen.\\
                $(+)$ = arriba, $(-)$ = abajo.
        \end{itemize}
      \end{definicionbox}
      \vspace{6pt}
      \begin{pasobox}[1]
        \AlmFontFootnote
        Para el punto $A(3,\,2)$:\\
        \quad$x = 3$: muévete 3 unidades a la \textbf{derecha}.\\
        \quad$y = 2$: muévete 2 unidades hacia \textbf{arriba}.
      \end{pasobox}
    \end{column}
    \begin{column}{0.46\textwidth}
      \begin{center}
        \begin{tikzpicture}[scale=0.75]
          % Ejes
          \draw[recto, ->] (-0.5,0) -- (4.5,0)
            node[right, font=\AlmFontSmall\bfseries, color=AlmAzulM] {$x$};
          \draw[recto, ->] (0,-0.5) -- (0,3.5)
            node[above, font=\AlmFontSmall\bfseries, color=AlmAzulM] {$y$};
          % Cuadrícula
          \foreach \i in {1,2,3,4}{
            \draw[AlmAzulC!25, line width=0.3pt] (\i,-0.4)--(\i,3.3);
          }
          \foreach \j in {1,2,3}{
            \draw[AlmAzulC!25, line width=0.3pt] (-0.4,\j)--(4.3,\j);
          }
          % Marcas
          \foreach \i in {1,2,3,4}{
            \node[font=\AlmFontTiny, color=AlmGris] at (\i,-0.25) {\i};
          }
          \foreach \j in {1,2,3}{
            \node[font=\AlmFontTiny, color=AlmGris] at (-0.22,\j) {\j};
          }
          \node[font=\AlmFontTiny, color=AlmGris] at (-0.22,-0.22) {0};
          % Punto A(3,2)
          \node[punto] (A) at (3,2) {};
          \node[font=\AlmFontSmall\bfseries, color=AlmAzulM, above right=1pt]
            at (A) {$A(3,\,2)$};
          % Flechas de desplazamiento
          \draw[AlmTerra, ->, line width=1.2pt, dashed]
            (0,0) -- (3,0)
            node[midway, below, font=\AlmFontTiny\bfseries, color=AlmTerra]
            {$x=3$};
          \draw[AlmVerde, ->, line width=1.2pt, dashed]
            (3,0) -- (3,2)
            node[midway, right, font=\AlmFontTiny\bfseries, color=AlmVerde]
            {$y=2$};
        \end{tikzpicture}
      \end{center}
    \end{column}
  \end{columns}
\end{frame}

% ---------------------------------------------------------------
\seccionframe{3. Los cuatro cuadrantes}
% ---------------------------------------------------------------

\begin{frame}{Los cuatro cuadrantes y sus signos}
  \begin{columns}[c]
    \begin{column}{0.44\textwidth}
      \begin{center}
        \begin{tikzpicture}[scale=0.80]
          % Zonas sombreadas
          \fill[AlmAzulC!22] (0,0) rectangle (2.8,2.8);   % I
          \fill[AlmAmbar!22] (-2.8,0) rectangle (0,2.8);  % II
          \fill[AlmTerra!18] (-2.8,-2.8) rectangle (0,0); % III
          \fill[AlmVerde!18] (0,-2.8) rectangle (2.8,0);  % IV
          % Ejes
          \draw[recto, ->] (-3.1,0) -- (3.2,0)
            node[right, font=\AlmFontSmall\bfseries, color=AlmAzulM] {$x$};
          \draw[recto, ->] (0,-3.1) -- (0,3.2)
            node[above, font=\AlmFontSmall\bfseries, color=AlmAzulM] {$y$};
          % Etiquetas de cuadrante
          \node[font=\AlmFontNormal\bfseries, color=AlmAzul]
            at (1.5, 1.5)  {I};
          \node[font=\AlmFontNormal\bfseries, color=AlmAmbar!80!black]
            at (-1.5, 1.5) {II};
          \node[font=\AlmFontNormal\bfseries, color=AlmTerra]
            at (-1.5,-1.5) {III};
          \node[font=\AlmFontNormal\bfseries, color=AlmVerde]
            at (1.5,-1.5)  {IV};
          % Signos en cuadrantes
          \node[font=\AlmFontTiny, color=AlmAzul]
            at (1.5, 0.9) {$(+,+)$};
          \node[font=\AlmFontTiny, color=AlmAmbar!80!black]
            at (-1.5, 0.9) {$(-,+)$};
          \node[font=\AlmFontTiny, color=AlmTerra]
            at (-1.5,-0.9) {$(-,-)$};
          \node[font=\AlmFontTiny, color=AlmVerde]
            at (1.5,-0.9)  {$(+,-)$};
          % Origen
          \node[font=\AlmFontTiny, color=AlmGris] at (-0.22,-0.22) {O};
        \end{tikzpicture}
      \end{center}
    \end{column}
    \begin{column}{0.52\textwidth}
      \vspace{4pt}
      {\AlmFontSmall
      \begin{tabular}{@{}cllc@{}}
        \toprule
        \textbf{Cuadrante} & \textbf{Signo $x$}
          & \textbf{Signo $y$} & \textbf{Color}\\
        \midrule
        \textbf{I}   & positivo $(+)$ & positivo $(+)$
          & \textcolor{AlmAzulM}{azul}\\
        \textbf{II}  & negativo $(-)$ & positivo $(+)$
          & \textcolor{AlmAmbar!70!black}{ámbar}\\
        \textbf{III} & negativo $(-)$ & negativo $(-)$
          & \textcolor{AlmTerra}{terracota}\\
        \textbf{IV}  & positivo $(+)$ & negativo $(-)$
          & \textcolor{AlmVerde}{verde}\\
        \bottomrule
      \end{tabular}}
      \vspace{8pt}
      \begin{notabox}[Truco para memorizar]
        \AlmFontSmall
        Empieza en el \textbf{cuadrante I} (arriba-derecha)
        y gira en sentido \textbf{contrario a las agujas del reloj}:
        I → II → III → IV.
      \end{notabox}
    \end{column}
  \end{columns}
\end{frame}

% ---------------------------------------------------------------
\seccionframe{4. Representar puntos en el plano}
% ---------------------------------------------------------------

\begin{frame}{Cinco puntos en el plano de Almería}
  \framesubtitle{Cada punto representa un lugar de la ciudad}
  \begin{columns}[c]
    \begin{column}{0.50\textwidth}
      \begin{center}
        \begin{tikzpicture}[scale=0.72]
          % Cuadrícula
          \foreach \i in {-4,-3,-2,-1,0,1,2,3,4,5}{
            \draw[AlmAzulC!20, line width=0.3pt]
              (\i,-4.3)--(\i,4.5);
          }
          \foreach \j in {-3,-2,-1,0,1,2,3,4}{
            \draw[AlmAzulC!20, line width=0.3pt]
              (-4.3,\j)--(5.3,\j);
          }
          % Ejes
          \draw[recto, ->] (-4.5,0) -- (5.5,0)
            node[right, font=\AlmFontSmall\bfseries, color=AlmAzulM] {$x$};
          \draw[recto, ->] (0,-3.8) -- (0,4.8)
            node[above, font=\AlmFontSmall\bfseries, color=AlmAzulM] {$y$};
          % Marcas numéricas
          \foreach \i in {-4,-3,-2,-1,1,2,3,4,5}{
            \node[font=\AlmFontTiny, color=AlmGris] at (\i,-0.3) {\i};
          }
          \foreach \j in {-3,-2,-1,1,2,3,4}{
            \node[font=\AlmFontTiny, color=AlmGris] at (-0.3,\j) {\j};
          }
          \node[font=\AlmFontTiny, color=AlmGris] at (-0.28,-0.28) {0};
          % Punto A(3,2)
          \node[punto] (A) at (3,2) {};
          \node[font=\AlmFontScript\bfseries, color=AlmAzulM,
                above right=0pt] at (A) {$A(3,2)$};
          % Punto B(-4,1)
          \node[puntot] (B) at (-4,1) {};
          \node[font=\AlmFontScript\bfseries, color=AlmTerra,
                above right=0pt] at (B) {$B(-4,1)$};
          % Punto C(-2,-3)
          \node[puntov] (C) at (-2,-3) {};
          \node[font=\AlmFontScript\bfseries, color=AlmVerde,
                above right=0pt] at (C) {$C(-2,-3)$};
          % Punto D(5,-1)
          \node[circle, fill=BCrear, inner sep=2.2pt,
                draw=AlmBlanco, line width=0.8pt] (D) at (5,-1) {};
          \node[font=\AlmFontScript\bfseries, color=BCrear,
                above left=0pt] at (D) {$D(5,-1)$};
          % Punto E(0,4)
          \node[circle, fill=BAplicar, inner sep=2.2pt,
                draw=AlmBlanco, line width=0.8pt] (E) at (0,4) {};
          \node[font=\AlmFontScript\bfseries, color=BAplicar!80!black,
                right=2pt] at (E) {$E(0,4)$};
        \end{tikzpicture}
      \end{center}
    \end{column}
    \begin{column}{0.46\textwidth}
      \vspace{4pt}
      {\AlmFontSmall
      \begin{tabular}{@{}clcc@{}}
        \toprule
        \textbf{Punto} & \textbf{Coordenadas}
          & \textbf{Cuadrante} & \textbf{Color}\\
        \midrule
        $A$ & $(3,\;2)$    & I   & \textcolor{AlmAzulM}{azul}\\
        $B$ & $(-4,\;1)$   & II  & \textcolor{AlmTerra}{rojo}\\
        $C$ & $(-2,\;-3)$  & III & \textcolor{AlmVerde}{verde}\\
        $D$ & $(5,\;-1)$   & IV  & \textcolor{BCrear}{púrpura}\\
        $E$ & $(0,\;4)$    & eje $y$ & \textcolor{BAplicar!80!black}{ámbar}\\
        \bottomrule
      \end{tabular}}
      \vspace{8pt}
      \begin{notabox}[Punto en el eje]
        \AlmFontFootnote
        Si $x = 0$, el punto está \textbf{sobre el eje Y}.\\
        Si $y = 0$, el punto está \textbf{sobre el eje X}.\\
        No pertenece a ningún cuadrante.
      \end{notabox}
    \end{column}
  \end{columns}
\end{frame}

% ---------------------------------------------------------------
\seccionframe{5. Escala: de la cuadrícula al mapa real}
% ---------------------------------------------------------------

\begin{frame}{Coordenadas y escala en el mapa de Almería}
  \begin{columns}[c]
    \begin{column}{0.52\textwidth}
      \begin{teoriabox}
        \AlmFontSmall
        En nuestro \textbf{mapa de Almería},
        asignamos la escala:
        \[
          1\ \text{unidad} = 100\ \text{m}
        \]
        Así, el punto $P(3,\,5)$ significa:
        \begin{itemize}
          \item $x=3$: 300 m hacia el \textbf{este}.
          \item $y=5$: 500 m hacia el \textbf{norte}.
        \end{itemize}
      \end{teoriabox}
      \vspace{6pt}
      \begin{ejemplobox}
        \AlmFontFootnote
        La \textbf{Alcazaba} de Almería tiene
        coordenadas GPS reales:\\
        \[
          36{,}8409°\,\text{N}, \quad 2{,}4699°\,\text{O}
        \]
        Estas coordenadas son \textbf{latitud/longitud},
        un sistema diferente al $(x,y)$ cartesiano.
      \end{ejemplobox}
    \end{column}
    \begin{column}{0.44\textwidth}
      \begin{center}
        \begin{tikzpicture}[scale=0.82]
          % Mini mapa con escala
          \fill[AlmAzulC!15, rounded corners=3pt]
            (0,0) rectangle (4.2,3.5);
          % Cuadrícula
          \foreach \i in {1,2,3,4}{
            \draw[AlmAzulC!40, line width=0.4pt] (\i,0)--(\i,3.5);
          }
          \foreach \j in {1,2,3}{
            \draw[AlmAzulC!40, line width=0.4pt] (0,\j)--(4.2,\j);
          }
          % Ejes
          \draw[recto, ->] (0,0)--(4.5,0)
            node[right,font=\AlmFontTiny\bfseries, color=AlmAzulM]{$x$};
          \draw[recto, ->] (0,0)--(0,3.8)
            node[above,font=\AlmFontTiny\bfseries, color=AlmAzulM]{$y$};
          % Marcas
          \foreach \i in {1,2,3,4}{
            \node[font=\AlmFontTiny, color=AlmGris] at (\i,-0.22) {\i};
          }
          \foreach \j in {1,2,3}{
            \node[font=\AlmFontTiny, color=AlmGris] at (-0.22,\j) {\j};
          }
          \node[font=\AlmFontTiny, color=AlmGris] at (-0.18,-0.18){0};
          % Punto P(3,2) — Alcazaba aprox.
          \node[puntot] (P) at (3,2) {};
          \node[font=\AlmFontTiny\bfseries, color=AlmTerra, above right=0pt]
            at (P) {Alcazaba};
          \node[font=\AlmFontTiny, color=AlmTerra, below right=0pt]
            at (P) {$(3,2)$};
          % Acotación x
          \draw[acota] (0,-0.48) -- (3,-0.48);
          \node[font=\AlmFontTiny, color=AlmAzulM] at (1.5,-0.72) {$300\ \text{m}$};
          % Acotación y
          \draw[acota] (3.5,0) -- (3.5,2);
          \node[font=\AlmFontTiny, color=AlmAzulM, rotate=90, anchor=south]
            at (3.75,1) {$200\ \text{m}$};
          % Escala
          \node[font=\AlmFontTiny, color=AlmGris, align=center]
            at (2.1,3.3) {escala: 1 ud = 100 m};
        \end{tikzpicture}
      \end{center}
    \end{column}
  \end{columns}
\end{frame}

% ---------------------------------------------------------------
\seccionframe{6. Ejercicios multinivel — Bloom}
% ---------------------------------------------------------------

\begin{frame}{Ejercicio COMPRENDER}
  \bloomtag{COMPRENDER}\quad\dificultad{2}
  \vspace{4pt}
  \begin{comprenderbox}
    \AlmFontSmall
    Responde razonando cada respuesta:
    \begin{enumerate}
      \item ¿Qué significa el par $(-3,\,4)$?
            Explica el movimiento en el plano:
            ¿hacia dónde y cuánto nos desplazamos en cada eje?
      \item ¿Por qué el par $(3,\,5)$ es diferente al par $(5,\,3)$?
            Representa \textbf{ambos} en un mismo sistema de ejes
            y comprueba que son puntos distintos.
      \item ¿Qué pares de coordenadas $(x, y)$ con $x = y$
            están a la misma distancia del origen
            sobre la recta diagonal?
            Escribe tres ejemplos.
    \end{enumerate}
  \end{comprenderbox}
\end{frame}

\begin{frame}{Ejercicio APLICAR}
  \bloomtag{APLICAR}\quad\dificultad{3}
  \vspace{4pt}
  \begin{aplicarbox}
    \AlmFontFootnote
    \textbf{Parte A — Representa los puntos:}\\
    Dibuja un sistema de ejes de $-2$ a $6$ en $x$
    y de $0$ a $6$ en $y$. Sitúa los puntos:
    \[
      A(1,\,4), \quad B(3,\,2), \quad C(5,\,5), \quad D(0,\,3)
    \]
    Une los puntos en orden $A \to B \to C \to D \to A$.
    ¿Qué forma geométrica se obtiene? Nómbrala.\\[6pt]
    \textbf{Parte B — Lee coordenadas del mapa:}\\
    Tu profesor proyectará un mapa con 4 puntos marcados.
    Escribe las coordenadas $(x, y)$ de cada uno.
    Compara con tus compañeros. ¿Todos obtuvisteis
    los mismos resultados?\\[6pt]
    \textbf{Parte C — Aplicación a Almería:}\\
    Si el origen $(0,0)$ es la Plaza de la Catedral y
    la escala es 1 ud = 100 m, escribe las coordenadas
    aproximadas de un monumento que conozcas.
  \end{aplicarbox}
\end{frame}

\begin{frame}{Ejercicio EVALUAR}
  \bloomtag{EVALUAR}\quad\dificultad{4}
  \vspace{4pt}
  \begin{evaluarbox}
    \AlmFontSmall
    \textbf{Debate y argumenta:}\\[6pt]
    El Ayuntamiento de Almería quiere crear una nueva
    guía turística digital.\\[4pt]
    \begin{itemize}
      \item ¿Deberían usar \textbf{coordenadas cartesianas}
            $(x, y)$ o \textbf{coordenadas GPS} (latitud/longitud)?
    \end{itemize}
    \vspace{6pt}
    Escribe un argumento de \textbf{5--6 líneas} que incluya:
    \begin{enumerate}
      \item Ventajas e inconvenientes de cada sistema.
      \item Tu recomendación justificada.
      \item Al menos un ejemplo concreto de Almería.
    \end{enumerate}
  \end{evaluarbox}
  \vspace{4pt}
  \begin{notabox}[Pista]
    \AlmFontFootnote
    Piensa en el público de la guía: ¿turistas con smartphone,
    escolares con papel, o profesionales de cartografía?
  \end{notabox}
\end{frame}

% ---------------------------------------------------------------
%  EVIDENCIA
% ---------------------------------------------------------------
\begin{frame}{Evidencia del día}
  \begin{evidenciabox}
    Entrega tu \textbf{Ficha de Coordenadas} con:
    \begin{itemize}
      \item Sistema de ejes claramente dibujado con regla,
            etiquetado $x$ e $y$.
      \item Al menos \textbf{6 puntos} correctamente colocados
            en las intersecciones de la cuadrícula
            (no entre líneas).
      \item Cada punto identificado con su \textbf{nombre}
            y su par $(x, y)$.
      \item El \textbf{cuadrante} al que pertenece indicado
            para cada punto.
      \item Una nota explicando con tus palabras
            la diferencia entre $(x, y)$ y GPS.
    \end{itemize}
  \end{evidenciabox}
\end{frame}

% ---------------------------------------------------------------
%  ERROR FRECUENTE
% ---------------------------------------------------------------
\begin{frame}{¡Cuidado con este error!}
  \begin{errorbox}
    \begin{columns}[c]
      \begin{column}{0.50\textwidth}
        \incorrecto\ \textbf{Error frecuente}\\[6pt]
        Confundir la coordenada $x$ con la $y$:
        \begin{center}
          \AlmFontLarge
          $B(1,\,-4) \;\neq\; B(-4,\,1)$
        \end{center}
        \AlmFontFootnote
        Si marcas el punto $B(-4,\,1)$ en la posición
        de $B(1,\,-4)$, estarás a \textbf{cientos de metros}
        del lugar real en el mapa de Almería.
      \end{column}
      \begin{column}{0.46\textwidth}
        \correcto\ \textbf{Correcto}\\[6pt]
        \AlmFontSmall
        \textbf{Recuerda el orden estricto:}
        \begin{enumerate}
          \item \textbf{Primera} $\to$ eje $x$ (horizontal).
          \item \textbf{Segunda} $\to$ eje $y$ (vertical).
        \end{enumerate}
        \vspace{4pt}
        \AlmFontFootnote
        $B(-4,\,1)$:\\
        \quad$x = -4$: 4 unidades a la \textbf{izquierda}.\\
        \quad$y = 1$: 1 unidad hacia \textbf{arriba}.
      \end{column}
    \end{columns}
  \end{errorbox}
  \vspace{6pt}
  \begin{notabox}[Regla de oro]
    \AlmFontSmall
    «Primero la \textbf{calle} ($x$ horizontal),
    luego el \textbf{piso} ($y$ vertical).»
  \end{notabox}
\end{frame}

% ---------------------------------------------------------------
%  RESUMEN
% ---------------------------------------------------------------
\begin{frame}{Resumen de la sesión 04}
  \begin{resumenbox}
    \begin{itemize}
      \item El sistema cartesiano tiene un \textbf{eje X}
            (horizontal) y un \textbf{eje Y} (vertical)
            que se cruzan en el \textbf{origen} $(0,0)$.
      \item Un par $(x,\,y)$: primero la coordenada
            \textbf{horizontal}, luego la \textbf{vertical}.
      \item Los cuatro cuadrantes tienen signos:
            I $(+,+)$, II $(-,+)$, III $(-,-)$, IV $(+,-)$.
      \item $(3,\,5) \neq (5,\,3)$:
            el \textbf{orden importa}.
      \item Con escala $1\ \text{ud} = 100\ \text{m}$
            podemos usar $(x,y)$ para el mapa de Almería.
    \end{itemize}
  \end{resumenbox}
  \vspace{6pt}
  \begin{center}
    {\AlmFontSmall\color{AlmAzulM}
    \faArrowCircleRight\enskip
    \textbf{Próxima sesión:}
    Rectas, segmentos y ángulos en la ciudad.}
  \end{center}
\end{frame}

\end{document}



